\documentclass[11pt]{article}
\input{packages.tex}
\usepackage{soul}
\usepackage{multicol}
\DeclareMathOperator{\sign}{sign}
\DeclareMathOperator{\Imm}{Im}
\DeclareMathOperator{\rg}{rg}

% Настройка колонтитулов
\pagestyle{fancy}
\fancyhf{} % Очистить все поля
\fancyhead[L]{Сделаем вид, что это не копипаст презентации :)}
\fancyhead[R]{\href{https://t.me/lost_luna}{Автор}}
\fancyfoot[C]{\thepage}               % Справа внизу

\setlength{\headheight}{13.6pt}

\hypersetup{
    colorlinks=true,
    linkcolor=blue,
    citecolor=blue,
    filecolor=blue,
    urlcolor=blue,
}

\begin{document}

\begin{titlepage}
    \centering
    \vspace*{\fill}  % заполняет пространство сверху

    {\Huge Лекция 4. Метод резолюций и краткое содержание предыдущих серий.
    \par}

    \vspace*{\fill}  % заполняет пространство снизу
\end{titlepage}

\setcounter{page}{1}

\section{Краткое содержание предыдущих серий}

\begin{definition}
    Формальная система -- это набор из четырёх множеств $\{\Sigma, \mathcal{F}, 
    \mathcal{A}, \mathcal{R}\}$, где \begin{itemize}
        \item $\Sigma$ -- алфавит (счётное множество);
        \item $\mathcal{F}$ -- формулы ($\mathcal{F} \subseteq \Sigma^*$ -- все 
        конечные слова в алфавите $\Sigma$);
        \item $\mathcal{A}$ -- аксиомы ($\mathcal{A} \subseteq \mathcal{F}$);
        \item $\mathcal{R}$ -- правила вывода (о них чуть поподробнее)
    \end{itemize}
\end{definition}

\subsection{Исчисление высказываний}

\subsubsection*{Алфавит}

\begin{itemize}
    \item Счетное множество $\mathcal{L}$ пропозициональных символов, которые 
    называют переменными
    \item Пропозициональные связки $\rightarrow$, $\lnot$ 
    \item Символы изменения приоритета: $"("$, $")"$
\end{itemize}

\subsubsection*{Формула}

\begin{definition}
    Слово $W$ в алфавите исчисления высказываний является формулой $\Longleftrightarrow$
    выполнено одно из трёх условий:
    \begin{itemize}
        \item $W$ --- переменная
        \item $W = (\lnot A)$, где $A$ --- формула
        \item $W = (A \to B)$, где $A$ и $B$ --- формулы
    \end{itemize}
\end{definition}

\subsubsection*{Схемы аксиом}
Самих аксиом бесконечно много, но схем всего три, вот наши красавицы:

\begin{itemize}
    \item \[( A \to (B \to A))\]
    \item  \[\left( \left(A \to \left(B \to C\right) \right) \to \left( \left(A 
    \to B \right) \to \left(A \to C \right) \right) \right) \]
    \item \[\left( \left( \left( \lnot B \right) \to \left( \lnot A\right) \right)
    \to \left( \left( \left( \lnot B\right) \to A \right) \to B \right) \right)\]
\end{itemize}

\subsubsection*{Правило вывода}
modeus ponens:  \[ A, \ A \to B \ \vdash \ B\]

\subsection{Исчисление предикатов}

\subsubsection*{Алфавит}

\begin{itemize}
    \item $\mathcal{X}$ --- счетное множество символов переменных
    \item $\mathcal{P}$ --- счетное множество предикатных символов 
    \item $\mathcal{F}$ --- счетное множество функциональных символов 
    \item Пропозициональные связки $\rightarrow$, $\lnot$ 
    \item Символ запятой: $"${,}$"$
    \item Квантор всеобщности: $\forall$
    \item Символы изменения приоритета: $"("$, $")"$
\end{itemize}

\subsubsection*{Терм}
\begin{definition}
    Слово $t$ является \textbf{термом} $\Longleftrightarrow$ выполнено одно из двух условий:
    \begin{itemize}
        \item $t$ --- переменная или константа (0-арный функциональный символ)
        \item $t = f(t_1, \dots, t_n)$, где $t_1, \dots, t_n$ --- термы, а $f$ 
        из множества $n$-арных функциональных символов
    \end{itemize}
\end{definition}

\subsubsection*{Элементарная формула}
\begin{definition}
    Слово в алфавите исчисления предикатов является \textbf{элементарной формулой}
    $\Longleftrightarrow$ оно имеет вид $A(t_1, \dots, t_n)$, где $t_1, \dots, t_n$
    --- термы, а $A$ из множества $n$-арных предикатных символов. 
\end{definition}

\subsubsection*{Формула}
\begin{definition}
    Слова $A$ в алфавите исчисления предикатов является \textbf{формулой} $\Longleftrightarrow$
    выполнено одно из четырех условий:
    \begin{itemize}
        \item $A$ --- элементарная формула 
        \item $A = (\lnot B)$, где $B$ --- формула 
        \item $A = (C \to B)$, где $B$ и $C$ --- формулы 
        \item $A = (\forall x B)$, где $B$ --- формула, а $x$ --- переменная
    \end{itemize}
\end{definition}

\begin{remark}
    Формула $(\exists x A)$ сокращенно обозначает формулу \[(\lnot (\forall x (\lnot A)))\]
\end{remark}

\subsubsection*{Схемы аксиом}
Красавиц становится больше:
\begin{itemize}
    \item \[( A \to (B \to A))\]
    \item  \[\left( \left(A \to \left(B \to C\right) \right) \to \left( \left(A 
    \to B \right) \to \left(A \to C \right) \right) \right) \]
    \item \[\left( \left( \left( \lnot B \right) \to \left( \lnot A\right) \right)
    \to \left( \left( \left( \lnot B\right) \to A \right) \to B \right) \right)\]
    \item \[ \forall x A \to A (x := t), \text{ если терм } t \text{ \textbf{свободен} (без паники,
    потом объясним)}\] \[\text{для переменной } x \text{ в формуле } A\]
    \item \[ \forall x (A \to B) \to (A \to \forall x B ), \text{ если в формуле } A 
    \text{ нет свободных вхождений переменной } x\]
\end{itemize}

\subsubsection*{Правила вывода}
\begin{itemize}
    \item Modus ponens: если выводимы формулы $A$ и $A \to B$, то выводима формула $B$. \[ A, \ A \to B \ \vdash \ B\]
    \item Генерализация: если выводима формула $A$, то выводима формула $\forall x A$
    \[A \vdash \forall x A\]
\end{itemize}

\begin{definition}
    В формуле $\forall x A$ подформула $A$ называется \textbf{областью действия} квантора.
\end{definition}

\begin{definition}
    Вхождение переменной $x$ называется \textbf{связанным}, если $x$ входит в область 
    действия некоторого квантора по $x$ или в кванторный блок $\forall x$. \\ 
    Остальные вхождения называются \textbf{свободными}. 
\end{definition}

\begin{definition}
    Переменные, имеющие свободные вхождения в формулу, называются \textbf{параметрами} формулы. 
\end{definition}

\begin{remark}
    Через $A(x := t)$ будем обозначать формулу, которая получается из $A$ подстановкой
    терма $t$ вместо всех свободных вхождений переменной $x$.
\end{remark}

\begin{definition}
    Будем говорить, что терм $t$ \textbf{свободен для переменной} $x$ в формуле $A$, 
    если ни одно свободное вхождение $x$ в $A$ не лежит в области действия квантора по 
    перменной, входящей в терм $t$.
\end{definition}

\subsection{Выводимость}

\begin{definition}
    Формула $A$ называется \textbf{выводимой} в исчислении высказываний, если существует
    последовательность формул \[A_1, A_2, \dots, A_n = A\]
    для каждой $A_i$ выполнено одно из двух условий:
    \begin{itemize}
        \item $A_i$ --- аксиома
        \item найдутся такие $j, k < i$, что $A_k = A_j \to A_i$ 
    \end{itemize}
    (то есть ранее в последовательности формул встречались и формула $A_j$, и формула 
    $A_j \to A_i$)
\end{definition}

\begin{definition}
    Формула $A$ называется \textbf{выводимой} в исчислении высказываний из множества гипотез, 
    если существует последовательность формул \[A_1, A_2, \dots, A_n = A\]
    для каждой $A_i$ выполнено одно из двух условий:
    \begin{itemize}
        \item $A_i$ --- аксиома
        \item $A_i$ --- гипотеза
        \item найдутся такие $j, k < i$, что $A_k = A_j \to A_i$ 
    \end{itemize}
    (то есть ранее в последовательности формул встречались и формула $A_j$, и формула 
    $A_j \to A_i$)
\end{definition}

\begin{definition}
    Формула $A$ называется \textbf{выводимой} в исчислении предикатов, если существует
    последовательность формул \[A_1, A_2, \dots, A_n = A\]
    для каждой $A_i$ выполнено одно из двух условий:
    \begin{itemize}
        \item $A_i$ --- аксиома
        \item найдутся такие $j, k < i$, что $A_k = A_j \to A_i$ 
        \item найдется такое $k < i$, что $A_i = \forall x A_k$ (то есть выводится
        применением правила вывода)
    \end{itemize}
\end{definition}

\begin{definition}
    Формула $A$ называется \textbf{выводимой} в исчислении высказываний из множества гипотез, 
    если существует последовательность формул \[A_1, A_2, \dots, A_n = A\]
    для каждой $A_i$ выполнено одно из двух условий:
    \begin{itemize}
        \item $A_i$ --- аксиома
        \item $A_i$ --- гипотеза
        \item найдутся такие $j, k < i$, что $A_k = A_j \to A_i$ 
        \item  найдется такое $k < i$, что $A_i = \forall x A_k$ (то есть выводится
        применением правила вывода)
    \end{itemize}
    (то есть ранее в последовательности формул встречались и формула $A_j$, и формула 
    $A_j \to A_i$)
\end{definition}

\subsection{Семантика в исчислении предикатов}
\begin{definition}
    Сигнатура --- множество предикатных и функциональных символов с указанием арностей.
\end{definition}

\begin{definition}
    Если в формулу входят только предикатные и функциональные символы из
    сигнатуры $\sigma$, то такая формула называется формулой в сигнатуре $\sigma$.
\end{definition}

\begin{definition}
    Интерпретацией $M$ сигнатуры $\sigma$ назовём:
    \begin{itemize}
        \item Непустое множество $M$ --- область интерпретации
        \item Соответствие интерпретации, которое каждому $k$-арному функциональному
        символу $f$ из $\sigma$ сопоставляет функцию $[f]: M^k \to M$, отображающую $M^k$ в $M$
        \item Соответствие интерпретации, которое каждому $k$-арному предикатному символу
        $A$ из $\sigma$ сопоставляет $k$-арное отношение $[A]$ на $M$, то есть $A \subseteq M$
    \end{itemize}
\end{definition}

\section{Правила оценки}
Пусть каждой переменной присвоено значение из области интерпретации --- оценка переменной. \\
Представим формулы и термы как деревья:
\begin{itemize}
    \item оценка листа --- это оценка переменной, или интерпретация константы, которой помечен лист
    \item оценка вершины, помеченной $k$-арным функциональным символом $f$, равна 
    $[f](v_1, \dots, v_k)$, где $v_i$ --- оценки потомков данной вершины
    \item оценка вершины, помеченной $k$-арным предикатным символом $A$, равна
    $[A](v_1, \dots, v_k)$, где $v_i$ --- оценки потомков данной вершины
    \item оценка вершины, помеченной символом $\lnot$, является отрицанием оценки её
    потомка (у такой вершины ровно один потомок)
    \item оценка вершины, помеченной символом $\to$, равна импликации $v_1 \to v_2$
    оценок $v_1$, $v_2$ её потомков
    \item оценка вершины, помеченной символом $\forall x$, равна истине тогда и только тогда,
    когда оценка её потомка истинна при всех возможных оценках переменных, отличающихся от 
    данной лишь значением переменной $x$.
\end{itemize}

\begin{definition}
    Оценкой формулы/терма является оценка корня соответствующего дерева.
\end{definition}

\section{Интерпретация формулы и терма}

\begin{definition}
    Интерпретацией терма $t$ является функция $[t]$, которая отображает набор $(a_1, \dots, a_k)$
    значений переменных $x_1, \dots, x_k$, входящих в терм, в оценку этого терма 
    при оценке переменных $x_i = a_i$.
\end{definition}

\begin{definition}
    Интерпретацией $[A]$ формулы A назовём $k$-арный предикат, где $k$ --- число
    параметров $x_1, \dots, x_k$ формулы $A$. Этот предикат истинен при наборе 
    значений параметров $(a_1, \dots, a_k)$  тогда и только тогда, когда истинна 
    оценка формулы $A$ при оценке переменных $x_1 = a_1$.
\end{definition}

\begin{definition}
    Формула без свободных вхождений переменных называется замкнутой.
\end{definition}

\begin{remark}
    Оценка формулы $A$ зависит только от значений параметров формулы $A$.
\end{remark}

\begin{remark}    
    Замкнутая формула интерпретируется как 0-арный предикат (истина или ложь).
\end{remark}

\section{Замыкание и общезначимость}

\begin{definition}
    Замыкание формулы A с параметрами $x_1, \dots, x_n$ --- это формула $\forall x_1 
    \dots \forall x_n$
\end{definition}

\begin{definition}
    Формула ИП называется общезначимой, если её замыкание истинно в любой
    интерпретации.
\end{definition}

\begin{definition}
    Формула ИП называется невыполнимой, если замыкание её отрицания общезначимо. 
    В противном случае формула называется выполнимой.
\end{definition}

\begin{remark}
    Иногда общезначимость определяют иначе: формула ИП называется общезначимой, 
    если она истинна в любой интерпретации при любой оценке переменных.
\end{remark}

\begin{definition}
    Формулы $A$ и $B$ эквивалентны, если формула $A \equiv B$общезначима
\end{definition}

\begin{multicols}{2}
    \begin{center}
        \textbf{Свойства ИВ}
    \end{center}

    \begin{itemize}
        \item Корректность:\\ В исчислении высказываний выводимы только тавтологии.
        \item Полнота:\\ Если формула является тавтологией, то она выводима в ИВ.
        \item Условная выводимость:\\ Формула семантически следует из системы гипотез тогда и только тогда, когда синтаксически следует (выводима) из этой системы.
    \end{itemize}

    \begin{center}
        \textbf{Свойства ИП}
    \end{center}

    \begin{itemize}
        \item Корректность:\\ В исчислении предикатов выводимы только общезначимые.
        \item Полнота:\\ Если формула является общезначимой, то она выводима в ИП.
        \item Условная выводимость:\\ Формула семантически следует из системы гипотез тогда и только тогда, когда синтаксически следует (выводима) из этой системы.
    \end{itemize}
\end{multicols}

\section{Метод резолюций}

\textbf{Сокращения:}
\begin{itemize}
    \item \[(A \wedge B) = \lnot (A \to (\lnot B))\]
    \item \[(A \vee B) = (\lnot A) \to B\]
\end{itemize}

\begin{definition}
    Литерал --- это переменная или отрицание переменной
\end{definition}

\begin{definition}
    Дизъюнкт --- это дизъюнкция литералов
\end{definition}

\begin{remark}
    Пустой дизъюнкт по определению означает ложь
\end{remark}

\begin{definition}
    Конъюнктивная нормальная форма (КНФ) --- конъюнкция дизъюнктов
\end{definition}

\begin{definition}
    Резольвентой дизъюнктов $\{x\} \cup A$ и $\{\lnot x\} \cup B$ относительно переменной
    $x$ называется дизъюнкт $A \cup B$.
\end{definition}

\begin{definition}
    Резолютивный вывод дизъюнкта $D$ из множества дизъюнктов $\mathcal{D}$ --- это такая
    последовательность дизъюнктов $D_1, \dots, D_n = D$, для каждого $D_i$ выполняется одно из двух условий:
    \begin{itemize}
        \item $D_i$ принадлежит $\mathcal{D}$
        \item найдутся такие $j, k < i$ что $D_i$ является резольвентой $D_j$ и $D_k$
    \end{itemize}
\end{definition}

\begin{theorem}[Корректность метода резолюций]
    Если существует вывод пустого дизъюнкта из дизъюнктов КНФ, то КНФ невыполнима.
\end{theorem}

\begin{definition}
    Назовем множество дизъюнктов $\mathcal{D}$ полным непротивочивым, если $\mathcal{D}$
    не содержит пустого дизъюнкта и любая резольвента дизъюнктов из $\mathcal{D}$
    либо тавтологична, либо содержится в $\mathcal{D}$.
\end{definition}

\begin{theorem}[Выполнимость]
    КНФ с полным непротиворечивым множеством дизъюнктов выполнима.
\end{theorem}

\section{Предварённая нормальная форма}

\begin{definition}
    Формула называется предварённой нормальной формой (ПНФ), если она имеет вид
    \[Q_1 x_1 \dots Q_k x_k A, \text{ где } Q_i \in \{\forall, \exists\}, \ A 
    \text{не содержит кванторов}\]
\end{definition}

\begin{remark}
    Для любой формулы существует эквивалентная ей предварённая нормальная форма.
\end{remark}

\begin{remark}
    Пусть $A, B$ — формулы исчисления предикатов, причём $B$ не содержит переменной
    $x$. Тогда следующие пары формул эквивалентны:
    \begin{itemize}
        \item \[ \lnot \exists x A \text{ и } \forall x \lnot A\]
        \item \[B \to \forall x A \text{ и } \forall x (B \to A)\]
        \item \[B \to \exists x A \text{ и } \exists x(B \to A)\]   
        \item \[\lnot \forall x A \text{ и } \exists x \lnot A\]   
        \item \[\forall x A \to B \text{ и } \exists x(A \to B)\]   
        \item \[\exists x A \to B \text{ и } \forall x (A \to B)\]
    \end{itemize}
\end{remark}

\begin{definition}
    Предварённая нормальная форма, которая содержит только кванторы
всеобщности, называется $\Pi_1$-формулой.
\end{definition}

\begin{itemize}
    \item Пусть $A$ — замкнутая формула в предварённой нормальной форме. \\
    Рассмотрим самый левый квантор существования $\exists y$, обозначим количество
    кванторов всеобщности перед ним через $k$.
    \item Введём новый функциональный символ $f$ арности $k$. Удалим из формулы
    квантор $\exists x$ и заменим в бескванторной части формулы $B$ все вхождения 
    $y$ на $f(x_1, \dots, x_k)$ или на $f$ при $k = 0$. Получим новую замкнутую формулу $C$.
    \item Будем повторять описанную выше процедуру до тех пор, пока не получим
    (замкнутую) $\Pi_1$-формулу
\end{itemize}

\section{Нормальная интерпретация}
\begin{itemize}
    \item Область нормальной интерпретации формул сигнатуры $\sigma$, содержащей
    константы, то есть 0-арные функциональные символы, состоит из замкнутых
    термов (строятся из констант) сигнатуры $\sigma$
    \item Если сигнатура $\sigma$ не содержит констант, то областью нормальной
    интерпретации будем полагать область интерпретации сигнатуры $\sigma \cap \{c\}$,
    где $c$ -- символ константы
    \item Нормальная интерпретация функциональных символов задаётся следующим 
    образом: применение функции $[f]$ переводит термы-объекты $\langle t_1 \rangle,
    \dots, \langle t_k \rangle$ в терм $\langle f(t_1, \dots, t_k) \rangle$.
    \item На интерпретацию предикатных символов ограничений нет
\end{itemize}

\begin{theorem}
    $\Pi_1$-формула выполнима $\Longleftrightarrow$ она выполнима 
    в нормальной интерпретации
\end{theorem}

\begin{theorem}[Эрбрана]
    Если замкнутая $\Pi_1$-формула $P$ с переменными $x_1, \dots, x_n$
    невыполнима, то существует такое конечное множество наборов
    замкнутых термов $t_1^{(i)}, \dots, t_n^{(i)}, \ 1 \leq i \leq k$, 
    что формула \[\bigwedge_{i = 1}^k F((x_1 := t_1^{(i)}), \dots, (x_n := t_n^{(i)}))\] 
    невыполнима. Здесь $F$ --- бескванторная часть формулы $P$.
\end{theorem}

Теперь распишем наш чудесный метод резолюций:
\begin{itemize}
    \item Строим ПНФ 
    \item Строим $\Pi_1$-форму
    \item Выражаем бескванторную часть в виде КНФ 
    \item Фиксируем нумерацию замкнутых формул в сигнатуре 
    \item Выберем число $h$ и зафиксируем первые $h$ элементарных замкнутых формул 
    в выбранной нумерации. Построим КНФ $C(h)$ с такой подстановкой замкнутых термов 
    вместо переменных, что все возникающие элементарные замкнутые подформулы содержатся среди
    $h$ зафиксированных формул 
    \item Если из $C(h)$ выводится пустой дизъюнкт то формула невыполнима
\end{itemize}
\end{document}